\begin{beginnerstatement}[In simple terms: Tauri and the desktop layer]

Tauri uses the existing web frontend and displays it as a desktop application
in a webview, an embedded area for displaying web content. An operating system
such as Windows, macOS, or Linux controls the computer; here, Rust forms the
native application layer, which can interact with the operating system directly
and provide native functions. A native application is installed and runs on the
computer rather than exclusively in a conventional browser. Together with
permissions, a capability describes which native functions are allowed; the
Content Security Policy (CSP) additionally restricts the content and connections
that may be used. A sidecar is a separate helper program that can be distributed
with the application, but the baseline does not provide one for FastAPI. A
remote API, by contrast, runs on a remote server and is accessed over the
network. Packaging bundles the application into installable packages, while
signing cryptographically confirms their origin and integrity. Whether to use a
remote API, sidecar, or local commands, as well as packaging and signing,
remains a product decision here.

\end{beginnerstatement}
